{\huge\chapter{Постановка задачи и проблематика Vehicle ReID}}
\label{sec:Chapter1} \index{Chapter1}

\section{Формальная постановка}
Пусть имеется набор камер \(\mathcal{C}=\{c_i\}\), каждая из которых наблюдает сцену под своим ракурсом. Пусть \(\mathcal{X}\) — множество всех детекций транспортных средств, а каждой детекции \(x \in \mathcal{X}\) сопоставлены: идентификатор камеры \(cam(x)\in\mathcal{C}\) и (для обучения/оценки) идентификатор экземпляра автомобиля \(id(x)\in\mathbb{N}\).

Задача \emph{реидентификации} формулируется как поиск по изображению‑запросу \(q\) из одной камеры соответствующих изображений того же физического экземпляра в галерее \(\mathcal{G}\) из других камер: необходимо построить отображение \(f_\theta: \mathcal{X}\rightarrow\mathbb{R}^d\), переводящее изображение в эмбеддинг так, чтобы расстояния \(D(f_\theta(q), f_\theta(g))\) были малы для \(id(g)=id(q)\) и велики для иных \(g\). Традиционно \(\theta\) обучается на тренировочном множестве \(\mathcal{X}_{train}\), причём \emph{множество идентичностей на тесте не пересекается с обучающим} (\(ID_{train}\cap ID_{test}=\varnothing\)).

\section{Ограничения и особенности данных}
В работе \textbf{не используется распознавание номерных знаков}; сопоставление выполняется исключительно по визуальным признакам кузова и деталей. Ключевые сложности:
\begin{itemize}
  \item \textbf{Смена ракурса и масштаба}: вид спереди/сзади/сбоку, разные фокусные расстояния;
  \item \textbf{Вариативность освещения}: день/ночь, тени, блики, погодные условия;
  \item \textbf{Окклюзии и фон}: частичное перекрытие людьми/авто, динамичный фон;
  \item \textbf{Низкое разрешение}: удалённые объекты и компрессия видеопотока;
  \item \textbf{Доменный сдвиг между камерами}: цветопередача и экспозиция зависят от устройства и настроек.
\end{itemize}

\section{Метрики качества}
Оценивание проводится в постановке поиска/ранжирования.
\paragraph{CMC/Rank‑$k$.} Кривая кумулятивного совпадения (CMC) отражает вероятность того, что верный экземпляр находится в топ‑$k$ результата. Для одного запроса \(q\) CMC@k равна индикатору попадания истинного соответствия в первые \(k\) позиций.
\paragraph{Average Precision и mAP.} Для запроса \(q\) со множеством релевантных элементов \(\mathcal{R}_q\) \emph{Average Precision} вычисляется как усреднение прецизий на позициях релевантов:
\begin{equation}
AP(q) = \frac{1}{|\mathcal{R}_q|}\sum_{n=1}^{N} P@n \cdot \mathbf{1}\{r_n \in \mathcal{R}_q\},
\end{equation}
где \(r_n\) — элемент на позиции \(n\) в отсортированной галерее, \(P@n\) — прецизия на позиции \(n\). \emph{mAP} определяется как среднее \(AP(q)\) по всем запросам.

\section{Протоколы оценки и сплиты}
Для корректного сравнения используются общепринятые протоколы набора VeRi‑776: фиксированные сплиты train/query/gallery, исключение \emph{same‑camera} изображений из ранжирования (\emph{no-same-camera}), отчёт по CMC@1/5 и mAP. % \cite{VeRi2016}

\section{Проблематика и риски}
\begin{itemize}
  \item \textbf{Переобучение на камеру:} модель может подстраиваться под статистику конкретных камер. Решение: BNNeck/InstanceNorm, цветовые аугментации, смешивание доменов.
  \item \textbf{Шум в разметке:} ошибочные ID/кропы ухудшают обучение Triplet Loss. Решение: жёсткие правила PK‑семплинга, hard‑example mining с порогами.
  \item \textbf{Open‑set сценарий:} объект может отсутствовать в галерее; требуется пороговое решение «нет совпадения». Решение: калибровка расстояний/оценок схожести.
  \item \textbf{Скорость инференса:} поиск по большой галерее. Решение: off‑line индексация (FAISS), сжатие эмбеддингов, аппроксимирующий поиск.
\end{itemize}
